% !TEX root = ../main.tex
\chapter{随伴：自由構成と忘却}
\label{chap:ch37_adjunction}
\BookIndex{ずいはん}{随伴}
\BookIndex{じゆうものいど}{自由モノイド}
\BookIndex{ぼうきゃくかんしゅ}{忘却関手}
\BookIndex{foldMap}{\texttt{foldMap}}

\begin{chapterabstract}
本章では、圏論の発展的な概念である随伴を、ソフトウェアエンジニアが利用できる直感に絞って導入します。
中心となるのは、「自由に作る」操作と「構造を忘れる」操作の対応です。
型 \texttt{A} からリスト \texttt{List A} を作ると、要素を順番に並べるための最小限の結合構造が自動的に得られます。
一方、モノイドを単なる型として見るときには、結合演算や単位元という構造を忘れます。
本章では、自由構成と構造を忘れる操作が一対として自然に対応する理由を、free monoid と list を通して説明します。
\end{chapterabstract}

\begin{learninggoals}
\item 随伴を、まず「自由構成」と「忘却」の対応として説明できるようになります。
\item \texttt{List A} が free monoid の具体例であることを、\texttt{foldMap} の一意性として読めるようになります。
\item 設計上の「最小構造を自動生成する」という考え方を、DSL、設定、ログ集約に結び付けられるようになります。
\item Galois 接続を、前順序における随伴の例として位置付けられるようになります。
\item 本章のミニモデルと Mathlib の一般的な随伴の対応を説明できるようになります。
\end{learninggoals}

\section{自由構成と忘却の対応}
\label{sec:ch37-why}

前章までに、型、関数、同型、関手、自然変換、モノイド、モナド、Pullback、Galois 接続を見てきました。
これらは、ソフトウェアでよく現れる「変換」「合成」「保存性」「互換性」を整理するための道具です。

随伴は、これらより一段抽象的な概念です。
初めて見ると、定義に現れる hom 集合の自然同型や単位・余単位が難しく感じられるかもしれません。
しかし、随伴の入口として最も大切なのは、「一方の方向では構造を自由に作り、反対の方向では構造を忘れる」という直感です。

自由構成は、「最低限の約束だけを満たす構造を自動的に作る」操作です。
忘却は、「現在持っている構造のうち、注目しない部分を除外し、土台のデータだけを見る」操作です。
自由構成と忘却が対応するとき、その背後には随伴があります。

ソフトウェアに置き換えると、この考え方は珍しいものではありません。
たとえば、生のイベント列からログ集約用の構造を作ります。
文字列トークン列から DSL の構文木を作る場合や、設定項目の列から設定マージ用の構造を作る場合もあります。
逆に、構文木や設定オブジェクトから「含まれる値だけ」を取り出し、追加の規則をいったん忘れることもあります。
随伴は、このような作る操作と忘れる操作の関係を、非常に一般的に説明する言葉です。

\FigurePlaceholder{fig-ch37-free-forgetful-overview.png}{0.90}{自由構成と忘却}{fig:ch37-free-forgetful-overview}

図\ref{fig:ch37-free-forgetful-overview}の一方向では材料から構造を生成し、反対方向では構造を持つ対象から土台の型だけを取り出します。
随伴が主張するのは、後で述べる構造保存写像の対応であり、この二つの操作が互いに逆関数であるということではありません。

\section{「自由に作る」と「構造を忘れる」}
\label{sec:ch37-free-forgetful}

\subsection{構造を作る操作}

まず、集合または型 \texttt{A} があるとします。
この段階では、\texttt{A} の値をどのように結合するかは決まっていません。
たとえば、\texttt{A} がイベント型なら、イベントを一つずつ保持するだけであり、「イベント列を結合する」「空のイベント列を持つ」という構造はまだ明示されていません。

ここで \texttt{List A} を作ると、状況が変わります。
リストには空リスト \texttt{[]} があり、リスト同士を連結する \texttt{++} があります。
さらに、連結は結合律を満たし、空リストは単位元として振る舞います。
したがって、\texttt{List A} は自然にモノイドになります。

モノイドとは、型 \texttt{M}、単位元 \texttt{unit : M}、結合演算 \texttt{op : M -> M -> M} を持ち、次の三つの法則を満たす構造です。
\[
  \mathrm{op}(\mathrm{unit}, x) = x,
  \quad
  \mathrm{op}(x, \mathrm{unit}) = x,
  \quad
  \mathrm{op}(\mathrm{op}(x,y),z) = \mathrm{op}(x,\mathrm{op}(y,z)).
\]
リストでは、単位元が空リストであり、結合演算がリスト連結です。

\subsection{構造を忘れる操作}

反対に、モノイド \texttt{M} があるとします。
\texttt{M} には、単位元、結合演算、法則が備わっています。
しかし、場面によっては、その構造をいったん忘れ、単なる値の型として扱いたいことがあります。

たとえば、ログ集約用の \texttt{LogBuffer} がモノイド構造を持っていても、UI に一覧表示するときには中身の値だけに注目するかもしれません。
設定オブジェクトがマージ演算を持っていても、ある検査関数では「設定値の集合」としてだけ見るかもしれません。

忘却は、構造を破壊する操作ではありません。
その文脈で参照しない情報を、見ないことにする操作です。
型に付いている API、法則、制約、インスタンスを一時的に脇へ置き、土台の値だけを扱うと考えるとよいでしょう。

\subsection{随伴の最初の形}

自由構成と忘却が随伴になるとは、次の対応が自然に成り立つという意味です。

自由モノイドと忘却関手の随伴は、直感的には次の対応で表せます。
\[
  \{\texttt{List A} \to \texttt{M} \text{ のモノイド準同型}\}
  \cong
  \{\texttt{A} \to \texttt{M} \text{ の関数}\}.
\]
左辺は「リスト全体をモノイド \texttt{M} へ構造を保って写す方法」です。
右辺は「一つの要素を \texttt{M} へ写す方法」です。
リスト全体を写すモノイド準同型と、要素ごとの関数が一対一に対応することが、free monoid の普遍性です。

ここで重要なのは、\texttt{List A} から \texttt{M} への写像が、単なる関数ではなくモノイド準同型であることです。
すなわち、空リストを単位元へ送り、リスト連結を \texttt{M} の結合演算へ送らなければなりません。

この条件により、\texttt{List A} 全体での振る舞いは、各要素 \texttt{a : A} をどこへ送るかだけで完全に決まります。
その拡張を実装する関数が \texttt{foldMap} です。

\section{free monoid と list}
\label{sec:ch37-free-monoid-list}

\subsection{リストは「順番に並べる」最小構造}

型 \texttt{A} の値を使ってモノイドを作りたいとします。
最初から \texttt{A} の値同士を足したり、掛けたり、マージしたりする規則を決める必要はありません。
最も素朴な方法は、値を並べることです。

\texttt{[a1, a2, a3]} は、\texttt{a1}、\texttt{a2}、\texttt{a3} をこの順番で並べたものを表します。
二つのリストを連結すれば、並びをつなげられます。
空リストは、何も並べていない状態です。
この構造は、\texttt{A} に余計な演算を要求しません。

\texttt{List A} は、\texttt{A} の値を材料として、「空」と「連結」を持つモノイドを自由に生成します。
ここで「自由」とは、モノイド則から必要になる等式以外に、要素同士の同一視、順序交換、簡約を追加しないという意味です。
構造の個数や包含関係について「最小」という意味ではありません。

\FigurePlaceholder{fig-ch37-free-monoid-list.png}{0.90}{自由モノイドの普遍性}{fig:ch37-free-monoid-list}

\subsection{foldMap が普遍性を実装する}

モノイド \texttt{M} と関数 \texttt{f : A -> M} があるとします。
\texttt{f} は、\texttt{A} の要素を一つずつ \texttt{M} の値へ変換するだけです。
これをリスト全体へ拡張するには、次のように処理します。

\begin{itemize}
\item 空リストは、\texttt{M} の単位元へ送ります。
\item \texttt{x :: xs} は、\texttt{f x} と \texttt{xs} の変換結果を \texttt{M} の演算で結合します。
\end{itemize}

これが \texttt{foldMap} です。

本章では、Mathlib の \texttt{Monoid} や \texttt{CategoryTheory.Adjunction} には進まず、自作した小さな \texttt{TinyMonoid} を使って説明します。
目的は、随伴を完全に形式化することではなく、free monoid の普遍性を「要素への関数から、リスト全体の準同型が一意に決まる」という形で読むことです。

\begin{listing}[htbp]
\begin{minted}{lean4}
structure TinyMonoid (M : Type) where
  unit : M
  op : M -> M -> M
  unit_left : forall x : M, op unit x = x
  unit_right : forall x : M, op x unit = x
  assoc : forall x y z : M,
    op (op x y) z = op x (op y z)

def listMonoid (A : Type) : TinyMonoid (List A) where
  unit := []
  op := List.append
  unit_left := by
    intro xs
    rfl
  unit_right := by
    intro xs
    induction xs with
    | nil => rfl
    | cons x xs ih =>
        simp [List.append]
  assoc := by
    intro xs ys zs
    induction xs with
    | nil => rfl
    | cons x xs ih =>
        simp [List.append]

def foldMap {A M : Type} (mon : TinyMonoid M)
    (f : A -> M) : List A -> M
  | [] => mon.unit
  | x :: xs => mon.op (f x) (foldMap mon f xs)

theorem foldMap_singleton {A M : Type}
    (mon : TinyMonoid M) (f : A -> M) (a : A) :
    foldMap mon f [a] = f a := by
  simp [foldMap, mon.unit_right]
\end{minted}
\caption{\texttt{TinyMonoid}・\texttt{listMonoid}・\texttt{foldMap}}
\label{lst:ch37_tiny_monoid}
\end{listing}

コード\ref{lst:ch37_tiny_monoid}では、\texttt{TinyMonoid} に単位元、二項演算、三つの法則を持たせています。
さらに、\texttt{listMonoid} では、空リストと連結がその法則を満たすことも示しています。
\texttt{foldMap} は、空リストと \texttt{x :: xs} に場合分けして、リストを \texttt{M} の値へ畳み込みます。
\texttt{foldMap\_singleton} は、1要素のリストに対して元の関数 \texttt{f} と同じ働きをすることを示しています。

\subsection{モノイド準同型としての拡張}

\texttt{foldMap} が単なる便利な関数ではない理由は、リスト連結を保存するからです。
すなわち、\texttt{xs ++ ys} を一度に \texttt{foldMap} した場合と、\texttt{xs} と \texttt{ys} を別々に \texttt{foldMap} してから \texttt{M} 側で結合した場合とで、結果が一致します。

\begin{listing}[htbp]
\begin{minted}{lean4}
theorem foldMap_append {A M : Type}
    (mon : TinyMonoid M) (f : A -> M) :
    forall xs ys : List A,
      foldMap mon f (xs ++ ys) =
        mon.op (foldMap mon f xs) (foldMap mon f ys) := by
  intro xs
  induction xs with
  | nil =>
      intro ys
      simp [foldMap, mon.unit_left]
  | cons x xs ih =>
      intro ys
      simp [foldMap, ih, mon.assoc]

structure MonoidHom (M N : Type)
    (monM : TinyMonoid M) (monN : TinyMonoid N) where
  toFun : M -> N
  map_unit : toFun monM.unit = monN.unit
  map_op : forall x y : M,
    toFun (monM.op x y) = monN.op (toFun x) (toFun y)

def foldMapHom {A M : Type} (mon : TinyMonoid M) (f : A -> M) :
    MonoidHom (List A) M (listMonoid A) mon where
  toFun := foldMap mon f
  map_unit := rfl
  map_op := by
    intro xs ys
    exact foldMap_append mon f xs ys

theorem foldMapHom_unique {A M : Type} (mon : TinyMonoid M) (f : A -> M)
    (h : MonoidHom (List A) M (listMonoid A) mon)
    (h_single : forall a : A, h.toFun [a] = f a) :
    forall xs : List A, h.toFun xs = foldMap mon f xs := by
  intro xs
  induction xs with
  | nil =>
      simpa [foldMap] using h.map_unit
  | cons x xs ih =>
      calc
        h.toFun (x :: xs) = h.toFun ([x] ++ xs) := rfl
        _ = mon.op (h.toFun [x]) (h.toFun xs) := h.map_op [x] xs
        _ = mon.op (f x) (foldMap mon f xs) := by
          simp [h_single x, ih]
        _ = foldMap mon f (x :: xs) := rfl
\end{minted}
\caption{\texttt{foldMap\_append}・\texttt{foldMapHom\_unique}}
\label{lst:ch37_foldmap_hom}
\end{listing}

\texttt{foldMap\_append} は、\texttt{foldMap} がリスト連結を保存することを表しています。
\texttt{foldMapHom\_unique} は、同じ要素写像を持つ任意の準同型 \texttt{h} が、すべてのリストに対して \texttt{foldMap} と同じ値を返すことを示しています。
後者の結論は点ごとの関数等式であり、\texttt{MonoidHom} 構造体同士の等式を直接述べてはいませんが、ここで必要となる一意性を与えています。
証明ではリストの帰納法を使い、空リストの場合には単位元の保存を、\texttt{x :: xs} の場合には演算の保存と帰納法の仮定を使います。

ここまでの内容から、次の対応が見えてきます。

\begin{center}
\begin{tabular}{lll}
\toprule
視点 & 右辺のデータ & 左辺の構造保存写像 \\
\midrule
自由構成 & 要素ごとの関数 \texttt{A -> M} & \texttt{List A -> M} のモノイド準同型 \\
実装 & 1件変換 & バッチ全体の畳み込み \\
設計 & 材料の意味づけ & 構造を壊さない集約 \\
\bottomrule
\end{tabular}
\end{center}

\subsection{一意性が意味すること}

free monoid の普遍性では、\texttt{f : A -> M} を拡張するモノイド準同型は存在するだけでなく、一意です。
この一意性は実務でも重要です。

たとえば、一つのイベントをログ集約値に変換する規則 \texttt{f} を決めます。
空のイベント列を単位元として扱い、イベント列の連結をログ集約値の結合に対応させます。
この二つの条件を守る限り、イベント列全体を処理する方法は \texttt{foldMap f} 以外にありません。

一意性には、「実装の選択肢を減らす」という効果があります。
要素レベルの意味付けと、空および結合を保つという設計方針が決まれば、リスト全体の処理は自動的に決まります。
これは、バッチ処理、ログ集約、設定生成、DSL の評価器で役立つ見方です。

コード\ref{lst:ch37_foldmap_hom}の\texttt{foldMapHom\_unique}が、この一意性を示しています。
証明ではリストに対する帰納法を使い、空リストの場合には単位元の保存を使います。
\texttt{x :: xs} の場合には \texttt{[x] ++ xs} と見なし、準同型が結合を保存することを使います。

\section{随伴の定義を少しだけ見る}

\label{sec:ch37-adjunction-definition}

一般の随伴は、Hom 集合の全単射だけではなく、両変数について自然な同型として定義されます~\cite{leinster-basic-category-theory,riehl-category-theory-in-context}。

ここでは、一般の随伴の定義を少しだけ確認します。
二つの圏 \(\mathcal{C}\) と \(\mathcal{D}\) があり、関手
\[
  F : \mathcal{C} \to \mathcal{D},
  \qquad
  U : \mathcal{D} \to \mathcal{C}
\]
があるとします。
\(F\) が左随伴であり、\(U\) が右随伴であるとは、\(\mathcal{C}\) の対象 \(A\) と \(\mathcal{D}\) の対象 \(M\) について、次の対応が自然に存在することをいいます。
\[
  \mathrm{Hom}_{\mathcal{D}}(F A, M)
  \cong
  \mathrm{Hom}_{\mathcal{C}}(A, U M).
\]

free monoid の例では、\(\mathcal{C}\) は型と関数の世界であり、\(\mathcal{D}\) はモノイドとモノイド準同型の世界です。
\(F A\) は \texttt{List A} であり、\(U M\) はモノイド \texttt{M} から構造を忘れた土台の型です。

随伴 \(F \dashv U\) は、「\(F A\) から \(M\) へ構造を保って写すこと」と、「\(A\) から \(U M\) へ通常の関数として写すこと」が同じ情報を持つという対応です。
free monoid では、\texttt{List A} から \texttt{M} への準同型と、\texttt{A} から \texttt{M} への関数が対応します。

随伴は、二つの型が同型であるという主張ではありません。
\texttt{A} と \texttt{List A} が同じであるという意味でもありません。
随伴が比較しているのは、対象そのものではなく、「その対象から、またはその対象への写し方」です。
この点を混同すると、随伴が単なる双方向変換のように見えてしまいます。

\section{設計上の意味：最小の構造を自動生成する}
\label{sec:ch37-design}

\subsection{DSL の構文を自由に作る}

DSL を設計するとき、最初から複雑な意味論を持たせる必要はありません。
まず、コマンドを並べる、選択肢を作る、名前を束縛するといった構文を自由に作ります。
その後、別の世界へ解釈する関数を与えます。

リストの例では、材料 \texttt{A} から \texttt{List A} を作ります。
DSL では、トークンや命令の集合から構文木を作ります。
一般の構文木は自由モノイドではなく、演算記号とその引数の個数を定めたシグネチャに対する自由代数として扱います。
両者に共通するのは、「材料から、指定した演算以外の等式を勝手に加えずに構造を作り、後から意味を与える」という設計です。

\subsection{設定生成とログ集約}

設定項目の列から設定オブジェクトを作る場合にも、同じ見方ができます。
一つの設定項目をどのように解釈するかを決め、それらを結合する規則をモノイドとして用意すれば、設定項目のリスト全体を \texttt{foldMap} で処理できます。

ログ集約でも同様です。
一つのイベントをログ断片へ変換する関数と、ログ断片を結合するモノイドがあれば、イベント列全体のログは自然に決まります。

この見方の利点は、「集約処理を簡潔に書ける」ことだけではありません。
要素ごとの変換、空入力の扱い、結合の扱いが分離されるため、仕様レビューで確認すべき点が明確になります。
\texttt{foldMap} は、自由構成から生じる設計パターンとして読めます。

\subsection{実装と仕様の距離}

本章の Lean モデルは非常に小さなものです。
現実のログ集約や設定マージには、順序依存、重複キー、優先順位、エラー、外部状態などが含まれます。
それでも、核となる性質をモノイドとして切り出せるなら、\texttt{foldMap} に基づく設計を利用できます。

ここで重要なのは、「実装のすべてを自由モノイドとして見る」ことではありません。
重要なのは、自由構成の形が現れている箇所を見つけることです。

\section{Galois 接続との関係}
\label{sec:ch37-galois}

\ChapterRef{chap:relations-preorders-galois}では、Galois 接続を抽象化と具体化の対応として捉えました。
随伴を学ぶと、Galois 接続を特別な随伴として理解できます。

前順序は、圏として見ることができます。
対象は値であり、\(a \leq b\) が成り立つとき、\(a\) から \(b\) へただ一つの射があると考えます。
このとき、二つの前順序の間の単調写像 \(\alpha\) と \(\gamma\) が次の条件を満たせば、両者は随伴です。
\[
  \alpha(a) \leq b
  \quad\Longleftrightarrow\quad
  a \leq \gamma(b).
\]

\FigurePlaceholder{fig-ch37-galois-adjunction.png}{0.90}{Galois 接続と随伴}{fig:ch37-galois-adjunction}

図\ref{fig:ch37-galois-adjunction}は、左辺の比較と右辺の比較が同値になることを、前順序を圏と見たときの Hom の対応として表しています。

\begin{listing}[htbp]
\begin{minted}{lean4}
structure PreorderMini (A : Type) where
  le : A -> A -> Prop
  refl : forall a, le a a
  trans : forall a b c, le a b -> le b c -> le a c

structure GaloisConnectionMini {A B : Type}
    (ordA : PreorderMini A) (ordB : PreorderMini B)
    (lower : A -> B) (upper : B -> A) : Prop where
  condition : forall a b,
    ordB.le (lower a) b <-> ordA.le a (upper b)
\end{minted}
\caption{\texttt{PreorderMini}・\texttt{GaloisConnectionMini}}
\label{lst:ch37_galois}
\end{listing}

\texttt{GaloisConnectionMini} は、本書の説明用に定義した最小構造です。
Mathlib にある同名概念の API と自作構造を混同しないように、\texttt{Mini} を付けています。
この定義では、\ChapterRef{chap:relations-preorders-galois}と同様に、反射律と推移律を持つ \texttt{PreorderMini} を明示的に渡します。
表示した同値条件から \texttt{lower} と \texttt{upper} の単調性も導けるため、フィールドとして重ねて要求していません。

自由モノイドの随伴では、「構造を保つ写像」と「要素への写像」が対応しました。
Galois 接続では、「抽象側で比較できること」と「具体側で比較できること」が対応します。
どちらも、一方の世界における問いを、もう一方の世界における問いへ翻訳しています。

この関係は、静的解析、権限設計、公開ログへのマスキング、仕様弱化で役立ちます。
詳細な世界から抽象的な世界へ移るとき、何が保存され、何を戻せるのかを、等式ではなく順序で表します。
このときの対応が Galois 接続であり、圏論的には前順序上の随伴です。

\section{本章のまとめ}
\label{sec:ch37-summary}

随伴は、まず「自由に作る操作」と「構造を忘れる操作」の対応として理解するとよいでしょう。

\texttt{List A} は、型 \texttt{A} から自由モノイドを作る典型例です。

\texttt{foldMap} は、要素ごとの関数 \texttt{A -> M} を、リスト全体のモノイド準同型 \texttt{List A -> M} へ拡張します。

free monoid の普遍性は、「そのような拡張が一意に決まる」という設計上の強い制約を与えます。

Galois 接続は、前順序を圏として見たときの随伴であり、抽象化と具体化の対応を表します。

自作のミニ定義により、自由構成、忘却、普遍性、Galois 接続という随伴の核を確認しました。
